\section{Literature Review and Theoretical Framework}

Evaluating urban living conditions across expanding metropolitan regions in the developing world draws on four complementary fields: remote sensing of socioeconomic indicators, 3D building morphology, pedestrian accessibility modeling, and graph neural networks for spatial anomaly detection. We synthesize these areas into a four-quadrant divergence framework that contrasts asset-based wealth with in-situ physical living standards.

\subsection{Remote Sensing and Machine Learning for Socioeconomic Estimation}
Satellite Earth observation has long supported socioeconomic monitoring. Nocturnal luminosity (DMSP and VIIRS) serves as an established proxy for regional economic activity and electrification \citep{Henderson2012, Mellander2015}, but encounters saturation in commercial cores and blooming across unpopulated surfaces \citep{Zhao2019}. To improve spatial resolution, transfer learning models use daytime multi-spectral imagery to predict nighttime radiance, learning visual features that correlate with local survey expenditures \citep{Burke2021, Jean2016, Yeh2020}. Meta's Relative Wealth Index (RWI; \citealt{Chi2022}) scaled this approach to 135 countries at 2.4~km resolution by training gradient-boosted trees on satellite, mobility, and connectivity data against Demographic and Health Surveys (DHS). 

However, survey asset targets capture household durable goods ownership \citep{Rutstein2004}. In dense urban settlements, asset accumulation can decouple from physical living conditions: households may own mobile devices and appliances while residing in overcrowded structures with unpaved access and open drainage canals.

\subsection{Built-Environment Morphology and Living Space Sufficiency}
Urban remote sensing has traditionally identified informal settlements through distinct morphological signatures: dense building configurations, organic street layouts, small footprints, and heterogeneous roofing \citep{Kohli2012, Sliuzas2010}. Early studies framed detection as binary ``slum vs. formal'' image classification \citep{Kuffer2016, Kuffer2020, Wurm2019}. While computationally effective, binary classifications obscure progressive consolidation \citep{Dovey2020} and ignore the multidimensional continuum of living conditions \citep{Abascal2022, Lilford2019}.

To measure conditions continuously, UN-Habitat \citep{UNHabitat2003, UNHabitat2020} establishes five shelter deprivation criteria: durable construction, access to improved water, adequate sanitation, tenure security, and sufficient living area ($<9.0\text{ m}^2$ per person). Estimating floor area per capita historically proved difficult from 2D overhead imagery. Recent open planetary-scale building datasets, particularly Google Open Buildings v3 \citep{Sirko2021} with 2.5D heights, now enable volumetric morphology estimation across the Global South \citep{Biljecki2023, Biljecki2024}. Intersecting building volumes with ambient population data (WorldPop; \citealt{Tatem2017}) and functional residential masks (JRC GHS-BUILT-C; \citealt{Schiavina2022}) enables continuous estimation of living floor area per capita across regional extents.

\subsection{Pedestrian Accessibility and Open Spatial Data}
Physical living standards also depend on spatial access to basic public amenities, including healthcare, education, retail, and green space \citep{Geurs2004, Handy1997}. The ``15-minute city'' concept emphasizes reachable daily services within a 15-minute walk \citep{Moreno2021}. Cumulative opportunity thresholds (destinations within a fixed walking catchment at 4.0~km/h) offer intuitive, policy-relevant metrics compared to gravity models \citep{Hosford2022, Weiss2018}.

In fragmented cities, Euclidean buffer distances can substantially overestimate accessibility \citep{Boeing2017}. Rapid suburban expansion frequently situates enclosed master-planned developments beside informal settlements \citep{Dovey2021, Firman2004}. Perimeter walls and highway barriers require multi-kilometer pedestrian detours, making topological network routing essential. While OpenStreetMap enables open routing \citep{Boeing2017}, volunteer contributions remain uneven in peripheral areas \citep{BarringtonLeigh2017}. The Overture Maps Foundation \citep{Overture2023, Overture2024} addresses this gap by combining open registries and algorithmic extraction into a standardized Global Entity Reference System (GERS).

\subsection{Graph Neural Networks and Spatial Anomaly Detection}
Unsupervised disparity detection across continuous space requires structural learning methods beyond linear spatial autocorrelation statistics \citep{Anselin1995, Reardon2004}. Graph Neural Networks (GNNs), such as Graph Convolutional Networks (GCNs; \citealt{Kipf2016}) and Graph Attention Networks (GATs; \citealt{Velickovic2018}), learn representations over non-Euclidean spatial graphs where hexagonal cells form nodes and geometric adjacency forms edges \citep{Zhang2022}.

Deep GNNs often encounter ``over-smoothing,'' where repeated neighborhood message passing homogenizes node embeddings \citep{Chen2020, Kipf2016}. In unsupervised anomaly detection, this characteristic can be repurposed. By training a Graph Attention Autoencoder (GAT-AE) with a narrow bottleneck ($d=2$) over pure geometric adjacency, the model enforces local spatial continuity. Where an edge crosses an abrupt socio-spatial boundary, message passing aggregates conflicting feature distributions, resulting in elevated reconstruction errors at the interface. Drawing on attributed graph anomaly frameworks (DOMINANT; \citealt{Ding2019}), attributed error decomposition separates built-environment indicators from biophysical land-cover edges:
\begin{equation}
\text{Reconstruction Error}_i = \sum_{k=1}^M (x_{ik} - \hat{x}_{ik})^2
\end{equation}

\subsection{Theoretical Synthesis: The 4-Quadrant Divergence Framework}
We synthesize these domains into the \textbf{4-Quadrant Divergence Framework}, conceptualizing urban space along two axes: Consumer Asset Wealth (Meta RWI) and In-Situ Physical Living Standards (PLSI).

\begin{figure}[htbp]
\centering
\begin{verbatim}
                  High Physical Living Standard (PLSI)
                                  ▲
                                  │
                 Quadrant III     │     Quadrant I
            (Subsidized Formal /  │ (Integrated Affluent /
             Well-Serviced Low-   │  Master-Planned Enclaves)
             Income Peripheries)  │
                                  │
◄─────────────────────────────────┼─────────────────────────────────►
Low Asset Wealth (Meta RWI)       │       High Asset Wealth (Meta RWI)
                                  │
                 Quadrant II      │     Quadrant IV
            (Unserviced Informal  │  (Asset-Dense, Infrastructure-
             Deprivation / Rural  │   Deprived Informal Settlements)
             Poverty)             │
                                  ▼
                   Low Physical Living Standard (PLSI)
\end{verbatim}
\caption{The 4-Quadrant Divergence Framework contrasting asset-based wealth with in-situ physical living standards.}
\label{fig:four_quadrant_framework}
\end{figure}

\begin{itemize}
    \item \textbf{Quadrant I (High Asset, High Physical)}: Formal residential communities with large dwelling areas, tree cover, and walkable amenities.
    \item \textbf{Quadrant II (Low Asset, Low Physical)}: Unplanned settlements and impoverished peri-urban areas lacking private assets and public services.
    \item \textbf{Quadrant III (Low Asset, High Physical)}: Subsidized public housing with modest household assets but planned dwelling volumes, drainage, and paved access.
    \item \textbf{Quadrant IV (High Asset, Low Physical)}: Established dense informal settlements where households hold durable goods (appliances, vehicles) but endure floor-space deprivation ($<9.0\text{ m}^2$), open drainage, and elevated particulate air pollution.
\end{itemize}

Supervised models trained on survey asset indices often conflate Quadrant IV with Quadrant I, obscuring local infrastructure deficits. The unsupervised PLSI framework provides an empirical tool to evaluate both dimensions independently.
